\section{资料范围与考场速记}

\begin{warningbox}
本模板中的题目只用于演示结构。正式项目必须在此说明资料快照、收录范围、缺卷、截断、
回忆卷、勘误和答案来源，并在末尾给出逐题来源索引。
\end{warningbox}

\subsection{可直接使用的速记工具}

\begin{cheatbox}[title={紧致性与连续映射}]
\textbf{定义。}空间 \(X\) 紧致，是指 \(X\) 的每个开覆盖都有有限子覆盖。

\textbf{连续像。}若 \(X\) 紧致且 \(f:X\to Y\) 连续，则 \(f(X)\) 紧致。直接做法是把
\(f(X)\) 的开覆盖逐个取原像，得到 \(X\) 的开覆盖，再取有限子覆盖并推回去。

\textbf{闭子集。}紧致空间的闭子集紧致；Hausdorff 空间的紧致子集闭。

\textbf{判断顺序。}题目出现“任意开覆盖”时，先问能否通过原像、补集或有限交把它变成
已知紧致空间上的覆盖。详见
\knowledgelink{sec:compactness}{紧致性与连续性的自包含讲解}。
\end{cheatbox}

\subsection{\(N\) 天冲刺安排}

\noindent\begin{tabularx}{\textwidth}{@{}cX@{}}
\toprule
天数 & 任务 \\
\midrule
1--\(n_1\) & 学完基础定义、核心定理和完整证明，制作自己的错误清单。\\
\(n_1+1\)--\(n_2\) & 按题型训练；先口述识别线索和方法，再独立写规范答案。\\
\(n_2+1\)--\(n_3\) & 限时完成全部往年卷，按题目 ID 回填薄弱知识点。\\
\(n_3+1\)--\(N\) & 闭卷完成无答案模拟卷，只复习真实失分点与速记页。\\
\bottomrule
\end{tabularx}

\noindent 用户只需设置总天数 \(N\)；未指定时取 \(N=10\)。由 agent 根据资料量、
目标与可用时间设计 \(n_1,n_2,n_3\)，并满足 \(n_1<n_2<n_3<N\)。
